\QHeading{第1题　在互相平行的直线下面的（　　）里画“\(\bigcirc\)”。}

\par\noindent\makebox[\linewidth][c]{
\begin{tikzpicture}[x=1bp,y=-1bp,color=black,draw=black,text=black]
\begin{scope}\clip (56.0,131.0) rectangle (523.0,252.67);
% 高清原题重排：PDF 20，原书 15。不含解答。
\hdhead{15}{相交与平行——信息窗1}
\hdsection{1}{149}{在互相平行的直线下面的（　　）里画“\(\bigcirc\)”。}
\q{113}{204}{32}{13.5}{①}
\q{313}{204}{32}{13.5}{②}
\q{490}{204}{32}{13.5}{③}
\q{685}{204}{32}{13.5}{④}
\q{137}{314}{108}{13.5}{（　　）}
\q{337}{314}{108}{13.5}{（　　）}
\q{525}{314}{108}{13.5}{（　　）}
\q{708}{314}{108}{13.5}{（　　）}
\begin{scope}[scale=.6666667,draw=black,line width=.65pt]
\draw (145,228)--(212,238)--(145,255);
\draw (343,234)--(395,234) (343,261)--(395,261);
\draw (555,219)--(555,260) (522,260)--(587,260);
\draw (726,215)--(726,271) (757,215)--(757,271);
\end{scope}
\hdsection{2}{382}{下面每个图形中，哪两条线段互相垂直？填一填。}
\q{113}{437}{34}{13.5}{①}
\q{363}{437}{34}{13.5}{②}
\q{620}{437}{34}{13.5}{③}
\begin{scope}[scale=.6666667,draw=black,line width=.65pt]
\draw (155,465)--(277,465)--(216,526)--cycle;
\draw (426,458)--(505,458)--(530,485)--(549,531)--(426,531)--cycle;
\draw (661,537)--(728,445)--(797,537)--cycle;
\draw (728,445)--(728,537) (728,445)--(712,537);
\draw (144,672)--(307,672) (394,672)--(558,672) (650,673)--(810,673);
\end{scope}
\q{210}{443}{32}{13.5}{\(c\)}
\q{166}{492}{32}{13.5}{\(b\)}
\q{251}{490}{32}{13.5}{\(a\)}
\q{458}{436}{32}{13.5}{\(e\)}
\q{408}{483}{32}{13.5}{\(d\)}
\q{522}{457}{32}{13.5}{\(a\)}
\q{545}{497}{32}{13.5}{\(b\)}
\q{479}{534}{32}{13.5}{\(c\)}
\q{681}{465}{32}{13.5}{\(e\)}
\q{704}{487}{32}{13.5}{\(c\)}
\q{733}{489}{32}{13.5}{\(d\)}
\q{763}{472}{32}{13.5}{\(a\)}
\q{718}{541}{32}{13.5}{\(b\)}
\q{145}{596}{220}{13.5}{互相垂直的：}
\q{395}{596}{220}{13.5}{互相垂直的：}
\q{652}{596}{220}{13.5}{互相垂直的：}
\hdsection{3}{718}{判断下面各组直线的位置关系。（填序号。）}
\q{145}{774}{34}{13.5}{①}
\q{320}{774}{34}{13.5}{②}
\q{496}{775}{34}{13.5}{③}
\q{670}{774}{34}{13.5}{④}
\q{145}{878}{34}{13.5}{⑤}
\q{320}{878}{34}{13.5}{⑥}
\q{496}{878}{34}{13.5}{⑦}
\q{670}{878}{34}{13.5}{⑧}
\begin{scope}[scale=.6666667,draw=black,line width=.6pt]
\draw (194,783)--(223,831) (223,783)--(252,831);
\draw (398,776)--(398,839) (367,808)--(430,808);
\draw (563,779)--(589,837) (558,835)--(589,779);
\draw (709,790)--(789,790)--(713,831);
\draw (194,899)--(252,924) (194,924)--(252,899);
\draw (367,898)--(430,898) (367,925)--(430,925);
\draw (536,886)--(612,934) (557.2,936.6)--(590.8,883.4);
\draw (746,879)--(746,932) (716,932)--(778,932);
\foreach \x in {161,402,642}{
  \draw[rounded corners=10pt,dashed,draw=black!65] (\x,984) rectangle +(138,104);
  \draw (\x+113,1084)--(\x+137,1084)--(\x+131,1091)--(\x+118,1091)--cycle;
  \draw (\x+125,1057)--(\x+125,1084);
  \path[draw=black,fill=bookpink] (\x+123,1061)--(\x+113,1081)--(\x+123,1081)--cycle;
  \path[draw=black,fill=white] (\x+127,1067)--(\x+135,1081)--(\x+127,1081)--cycle;
}
\end{scope}
\q{204}{1130}{105}{14}{平行}
\q{443}{1130}{105}{14}{相交}
\q{680}{1130}{105}{14}{垂直}

\begin{scope}[color=answerred,draw=answerred,text=answerred]\draw[answerred,line width=.8pt] (246,214) circle[radius=6bp];
\draw[answerred,line width=.8pt] (493.3,214) circle[radius=6bp];
\ans{100.00}{435.330}{131.33}{11}{\(a\perp b\)}
\ans{265.33}{435.330}{130.67}{11}{\(c\perp d,\quad d\perp e\)}
\ans{439.33}{435.330}{112.00}{11}{\(b\perp d\)}
\draw[answerred,line width=.7pt] (147.38,345.95)--(142.67,341.24)--(137.95,345.95);
\draw[answerred,line width=.7pt] (284,348)--(290.67,348)--(290.67,354.67);
\draw[answerred,line width=.7pt] (284,312)--(290.67,312)--(290.67,305.33);
\draw[answerred,line width=.7pt] (485.33,351.33)--(492,351.33)--(492,358);
\ans{126.00}{683.33}{72.67}{12}{\(1,6\)}
\ans{273.33}{675.33}{90.00}{10}{\(2,3,4,5,7,8\)\\（包括垂直）}
\ans{445.33}{683.33}{71.33}{12}{\(2,7,8\)}
\ans{100.67}{773.33}{450.67}{9}{注：垂直也属于相交；若要求三栏不重复，相交栏填3、4、5。}
\end{scope}
\end{scope}\pgfresetboundingbox\path[use as bounding box] (56.0,131.0) rectangle (523.0,252.67);
\end{tikzpicture}}\par
\par\noindent




\par



\par\noindent

\QHeading{第2题　下面每个图形中，哪两条线段互相垂直？填一填。}

\par\noindent\makebox[\linewidth][c]{
\begin{tikzpicture}[x=1bp,y=-1bp,color=black,draw=black,text=black]
\begin{scope}\clip (56.0,285.0) rectangle (545.0,476.67);
% 高清原题重排：PDF 20，原书 15。不含解答。
\hdhead{15}{相交与平行——信息窗1}
\hdsection{1}{149}{在互相平行的直线下面的（　　）里画“\(\bigcirc\)”。}
\q{113}{204}{32}{13.5}{①}
\q{313}{204}{32}{13.5}{②}
\q{490}{204}{32}{13.5}{③}
\q{685}{204}{32}{13.5}{④}
\q{137}{314}{108}{13.5}{（　　）}
\q{337}{314}{108}{13.5}{（　　）}
\q{525}{314}{108}{13.5}{（　　）}
\q{708}{314}{108}{13.5}{（　　）}
\begin{scope}[scale=.6666667,draw=black,line width=.65pt]
\draw (145,228)--(212,238)--(145,255);
\draw (343,234)--(395,234) (343,261)--(395,261);
\draw (555,219)--(555,260) (522,260)--(587,260);
\draw (726,215)--(726,271) (757,215)--(757,271);
\end{scope}
\hdsection{2}{382}{下面每个图形中，哪两条线段互相垂直？填一填。}
\q{113}{437}{34}{13.5}{①}
\q{363}{437}{34}{13.5}{②}
\q{620}{437}{34}{13.5}{③}
\begin{scope}[scale=.6666667,draw=black,line width=.65pt]
\draw (155,465)--(277,465)--(216,526)--cycle;
\draw (426,458)--(505,458)--(530,485)--(549,531)--(426,531)--cycle;
\draw (661,537)--(728,445)--(797,537)--cycle;
\draw (728,445)--(728,537) (728,445)--(712,537);
\draw (144,672)--(307,672) (394,672)--(558,672) (650,673)--(810,673);
\end{scope}
\q{210}{443}{32}{13.5}{\(c\)}
\q{166}{492}{32}{13.5}{\(b\)}
\q{251}{490}{32}{13.5}{\(a\)}
\q{458}{436}{32}{13.5}{\(e\)}
\q{408}{483}{32}{13.5}{\(d\)}
\q{522}{457}{32}{13.5}{\(a\)}
\q{545}{497}{32}{13.5}{\(b\)}
\q{479}{534}{32}{13.5}{\(c\)}
\q{681}{465}{32}{13.5}{\(e\)}
\q{704}{487}{32}{13.5}{\(c\)}
\q{733}{489}{32}{13.5}{\(d\)}
\q{763}{472}{32}{13.5}{\(a\)}
\q{718}{541}{32}{13.5}{\(b\)}
\q{145}{596}{220}{13.5}{互相垂直的：}
\q{395}{596}{220}{13.5}{互相垂直的：}
\q{652}{596}{220}{13.5}{互相垂直的：}
\hdsection{3}{718}{判断下面各组直线的位置关系。（填序号。）}
\q{145}{774}{34}{13.5}{①}
\q{320}{774}{34}{13.5}{②}
\q{496}{775}{34}{13.5}{③}
\q{670}{774}{34}{13.5}{④}
\q{145}{878}{34}{13.5}{⑤}
\q{320}{878}{34}{13.5}{⑥}
\q{496}{878}{34}{13.5}{⑦}
\q{670}{878}{34}{13.5}{⑧}
\begin{scope}[scale=.6666667,draw=black,line width=.6pt]
\draw (194,783)--(223,831) (223,783)--(252,831);
\draw (398,776)--(398,839) (367,808)--(430,808);
\draw (563,779)--(589,837) (558,835)--(589,779);
\draw (709,790)--(789,790)--(713,831);
\draw (194,899)--(252,924) (194,924)--(252,899);
\draw (367,898)--(430,898) (367,925)--(430,925);
\draw (536,886)--(612,934) (557.2,936.6)--(590.8,883.4);
\draw (746,879)--(746,932) (716,932)--(778,932);
\foreach \x in {161,402,642}{
  \draw[rounded corners=10pt,dashed,draw=black!65] (\x,984) rectangle +(138,104);
  \draw (\x+113,1084)--(\x+137,1084)--(\x+131,1091)--(\x+118,1091)--cycle;
  \draw (\x+125,1057)--(\x+125,1084);
  \path[draw=black,fill=bookpink] (\x+123,1061)--(\x+113,1081)--(\x+123,1081)--cycle;
  \path[draw=black,fill=white] (\x+127,1067)--(\x+135,1081)--(\x+127,1081)--cycle;
}
\end{scope}
\q{204}{1130}{105}{14}{平行}
\q{443}{1130}{105}{14}{相交}
\q{680}{1130}{105}{14}{垂直}

\begin{scope}[color=answerred,draw=answerred,text=answerred]\draw[answerred,line width=.8pt] (246,214) circle[radius=6bp];
\draw[answerred,line width=.8pt] (493.3,214) circle[radius=6bp];
\ans{100.00}{435.330}{131.33}{11}{\(a\perp b\)}
\ans{265.33}{435.330}{130.67}{11}{\(c\perp d,\quad d\perp e\)}
\ans{439.33}{435.330}{112.00}{11}{\(b\perp d\)}
\draw[answerred,line width=.7pt] (147.38,345.95)--(142.67,341.24)--(137.95,345.95);
\draw[answerred,line width=.7pt] (284,348)--(290.67,348)--(290.67,354.67);
\draw[answerred,line width=.7pt] (284,312)--(290.67,312)--(290.67,305.33);
\draw[answerred,line width=.7pt] (485.33,351.33)--(492,351.33)--(492,358);
\ans{126.00}{683.33}{72.67}{12}{\(1,6\)}
\ans{273.33}{675.33}{90.00}{10}{\(2,3,4,5,7,8\)\\（包括垂直）}
\ans{445.33}{683.33}{71.33}{12}{\(2,7,8\)}
\ans{100.67}{773.33}{450.67}{9}{注：垂直也属于相交；若要求三栏不重复，相交栏填3、4、5。}
\end{scope}
\end{scope}\pgfresetboundingbox\path[use as bounding box] (56.0,285.0) rectangle (545.0,476.67);
\end{tikzpicture}}\par



\par





\par\noindent

\QHeading{第3题　判断下面各组直线的位置关系。（填序号。）}



\par


\par\noindent\makebox[\linewidth][c]{
\begin{tikzpicture}[x=1bp,y=-1bp,color=black,draw=black,text=black]
\begin{scope}\clip (91.0,511.0) rectangle (532.0,771.33);
% 高清原题重排：PDF 20，原书 15。不含解答。
\hdhead{15}{相交与平行——信息窗1}
\hdsection{1}{149}{在互相平行的直线下面的（　　）里画“\(\bigcirc\)”。}
\q{113}{204}{32}{13.5}{①}
\q{313}{204}{32}{13.5}{②}
\q{490}{204}{32}{13.5}{③}
\q{685}{204}{32}{13.5}{④}
\q{137}{314}{108}{13.5}{（　　）}
\q{337}{314}{108}{13.5}{（　　）}
\q{525}{314}{108}{13.5}{（　　）}
\q{708}{314}{108}{13.5}{（　　）}
\begin{scope}[scale=.6666667,draw=black,line width=.65pt]
\draw (145,228)--(212,238)--(145,255);
\draw (343,234)--(395,234) (343,261)--(395,261);
\draw (555,219)--(555,260) (522,260)--(587,260);
\draw (726,215)--(726,271) (757,215)--(757,271);
\end{scope}
\hdsection{2}{382}{下面每个图形中，哪两条线段互相垂直？填一填。}
\q{113}{437}{34}{13.5}{①}
\q{363}{437}{34}{13.5}{②}
\q{620}{437}{34}{13.5}{③}
\begin{scope}[scale=.6666667,draw=black,line width=.65pt]
\draw (155,465)--(277,465)--(216,526)--cycle;
\draw (426,458)--(505,458)--(530,485)--(549,531)--(426,531)--cycle;
\draw (661,537)--(728,445)--(797,537)--cycle;
\draw (728,445)--(728,537) (728,445)--(712,537);
\draw (144,672)--(307,672) (394,672)--(558,672) (650,673)--(810,673);
\end{scope}
\q{210}{443}{32}{13.5}{\(c\)}
\q{166}{492}{32}{13.5}{\(b\)}
\q{251}{490}{32}{13.5}{\(a\)}
\q{458}{436}{32}{13.5}{\(e\)}
\q{408}{483}{32}{13.5}{\(d\)}
\q{522}{457}{32}{13.5}{\(a\)}
\q{545}{497}{32}{13.5}{\(b\)}
\q{479}{534}{32}{13.5}{\(c\)}
\q{681}{465}{32}{13.5}{\(e\)}
\q{704}{487}{32}{13.5}{\(c\)}
\q{733}{489}{32}{13.5}{\(d\)}
\q{763}{472}{32}{13.5}{\(a\)}
\q{718}{541}{32}{13.5}{\(b\)}
\q{145}{596}{220}{13.5}{互相垂直的：}
\q{395}{596}{220}{13.5}{互相垂直的：}
\q{652}{596}{220}{13.5}{互相垂直的：}
\hdsection{3}{718}{判断下面各组直线的位置关系。（填序号。）}
\q{145}{774}{34}{13.5}{①}
\q{320}{774}{34}{13.5}{②}
\q{496}{775}{34}{13.5}{③}
\q{670}{774}{34}{13.5}{④}
\q{145}{878}{34}{13.5}{⑤}
\q{320}{878}{34}{13.5}{⑥}
\q{496}{878}{34}{13.5}{⑦}
\q{670}{878}{34}{13.5}{⑧}
\begin{scope}[scale=.6666667,draw=black,line width=.6pt]
\draw (194,783)--(223,831) (223,783)--(252,831);
\draw (398,776)--(398,839) (367,808)--(430,808);
\draw (563,779)--(589,837) (558,835)--(589,779);
\draw (709,790)--(789,790)--(713,831);
\draw (194,899)--(252,924) (194,924)--(252,899);
\draw (367,898)--(430,898) (367,925)--(430,925);
\draw (536,886)--(612,934) (557.2,936.6)--(590.8,883.4);
\draw (746,879)--(746,932) (716,932)--(778,932);
\foreach \x in {161,402,642}{
  \draw[rounded corners=10pt,dashed,draw=black!65] (\x,984) rectangle +(138,104);
  \draw (\x+113,1084)--(\x+137,1084)--(\x+131,1091)--(\x+118,1091)--cycle;
  \draw (\x+125,1057)--(\x+125,1084);
  \path[draw=black,fill=bookpink] (\x+123,1061)--(\x+113,1081)--(\x+123,1081)--cycle;
  \path[draw=black,fill=white] (\x+127,1067)--(\x+135,1081)--(\x+127,1081)--cycle;
}
\end{scope}
\q{204}{1130}{105}{14}{平行}
\q{443}{1130}{105}{14}{相交}
\q{680}{1130}{105}{14}{垂直}

\begin{scope}[color=answerred,draw=answerred,text=answerred]\draw[answerred,line width=.8pt] (246,214) circle[radius=6bp];
\draw[answerred,line width=.8pt] (493.3,214) circle[radius=6bp];
\ans{100.00}{435.330}{131.33}{11}{\(a\perp b\)}
\ans{265.33}{435.330}{130.67}{11}{\(c\perp d,\quad d\perp e\)}
\ans{439.33}{435.330}{112.00}{11}{\(b\perp d\)}
\draw[answerred,line width=.7pt] (147.38,345.95)--(142.67,341.24)--(137.95,345.95);
\draw[answerred,line width=.7pt] (284,348)--(290.67,348)--(290.67,354.67);
\draw[answerred,line width=.7pt] (284,312)--(290.67,312)--(290.67,305.33);
\draw[answerred,line width=.7pt] (485.33,351.33)--(492,351.33)--(492,358);
\ans{126.00}{683.33}{72.67}{12}{\(1,6\)}
\ans{273.33}{675.33}{90.00}{10}{\(2,3,4,5,7,8\)\\（包括垂直）}
\ans{445.33}{683.33}{71.33}{12}{\(2,7,8\)}

\end{scope}
\end{scope}\pgfresetboundingbox\path[use as bounding box] (91.0,511.0) rectangle (532.0,771.33);
\end{tikzpicture}}\par
\Answer{注：垂直也属于相交；若要求三栏不重复，相交栏填3、4、5。}


\QHeading{第1—3题　判定方法补充}

\Answer{方法一：延长线段所表示的直线，观察它们是否相交；不能因为纸上两小段暂时不碰到，就断定平行。方法二：用直尺和三角板检验方向；两直线相交成直角时互相垂直。垂直是相交的一种特殊情况，答案分类仍按本题原有分类口径核对。}
